\begin{abstract}
Every connected complete smooth Riemannian three-manifold without
boundary satisfying $-4\Delta+\Scal\ge\lambda>0$ as a quadratic-form
inequality admits a continuous map to a graph whose entire fibers have
diameter at most $500/\sqrt\lambda$ in the original metric.
No orientability, spin, compactness, or bounded-geometry assumption
is required.
\end{abstract}

\section{Introduction}\label{sec:introduction}

A positive scalar-curvature lower bound need not bound the diameter of a
three-manifold: the product of a round two-sphere with a line is a basic
example. A more suitable measure of size allows the manifold to extend
along a graph while requiring the inverse image of every point of the
graph to remain small. For a metric space $(X,d)$, its \emph{Urysohn
$1$-width} is
\[
 \operatorname{UW}_1(X,d)
 =\inf_{K,f}\ \sup_{y\in K}\diam_d f^{-1}(y),
\]
where $K$ ranges over simplicial complexes of dimension at most one, with
their usual polyhedral topology, and $f:X\to K$ ranges over continuous
maps. We give the empty fiber diameter zero. The diameter here is that of
the entire fiber, including distances between its different connected
components.

We use the sign convention $\Delta=\Div\nabla$ for the Laplacian.
The following theorem gives a width bound under a spectral, rather than
a pointwise, scalar-curvature hypothesis.

\begin{theorem}\label{thm:main}
Let $(M^3,g)$ be a connected complete smooth Riemannian three-manifold
without boundary. Suppose that $\lambda>0$ and
\begin{equation}\label{eq:main-spectral}
 \int_M\bigl(4|\nabla_g\phi|_g^2+\Scal_g\phi^2\bigr)\,d\vol_g
 \ge \lambda\int_M\phi^2\,d\vol_g
 \qquad\bigl(\phi\in C_c^\infty(M;\R)\bigr).
\end{equation}
Then there exist a simplicial complex $K$ of dimension at most one and a
continuous map $f:M\to K$ such that
\[
 \diam_g f^{-1}(y)\le \frac{500}{\sqrt\lambda}
 \qquad(y\in K).
\]
The complex has its usual polyhedral topology. No orientability, spin,
compactness, or bounded-geometry assumption is imposed on $M$.
\end{theorem}

The notation $-4\Delta_g+\Scal_g\ge\lambda$ will always mean
\eqref{eq:main-spectral}. In particular, a pointwise bound
$\Scal_g\ge\lambda$ implies the hypothesis, but is not required.
The theorem concerns the given metric $g$; passing to a conformally
related metric of positive scalar curvature would not by itself give
the required distance estimate. The constant $500$ is convenient and
is not optimized.

\subsection{Background and scope}

Stable minimal surfaces provide a classical link between scalar
curvature and the geometry of three-manifolds. Schoen and Yau used the
stability inequality and the Gauss equation to obtain curvature and
topological restrictions on such surfaces \cite{SchoenYau1979}.
Fischer-Colbrie and Schoen developed the associated positive-solution
method for Schr\"odinger operators \cite{FischerColbrieSchoen1980}.
The stable-surface radius estimate of Schoen and Yau
\cite{SchoenYau1983} and the distance-level estimates of Gromov and
Lawson \cite[Section~10]{GromovLawson1983} show how local surface
stability can control the large-scale geometry of the ambient
three-manifold. The latter distance-level argument uses a topological
condition on the first homology; it is not a diameter bound for an
arbitrary distance sphere in an arbitrary three-manifold.

Gromov's torus-stabilized scalar-curvature framework gives the closest
precedent for the spectral formulation here. Torus stabilization means
considering metrics on $M\times\mathbb T^N$ of the form
$g+\sum_{i=1}^N\varphi_i^2\,d\theta_i^2$, with smooth positive
functions $\varphi_i$ on $M$, while measuring width in the base metric
$g$. In his \emph{Four lectures on scalar curvature}, Gromov gives a
three-dimensional graph-width theorem in this framework
\cite[Section~3.10.1, Theorem~(G)]{GromovFourLectures2023}.
That statement is formulated for orientable complete three-manifolds,
with a mean-convex boundary allowed in the torus extension. The relation between torus
stabilization and the coefficient $4$ in
$-4\Delta+\Scal$ is made explicit in
\cite[Sections~1.A--1.B]{GromovStabilized2024}. Thus the spectral
viewpoint and the scalar-curvature-to-width principle are established
parts of this theory. Our purpose is to give a direct weighted proof
of Theorem~\ref{thm:main}, including its nonorientable and noncompact
cases, without a global bound on the auxiliary weight or its derivatives.

The prescribed-mean-curvature method used below also belongs to this
development. Gromov's $\mu$-bubbles minimize area with an added volume
term, whose density prescribes the mean curvature. Chodosh and Li
developed weighted versions and a slice-and-dice construction that first
simplifies the topology and then produces pieces of bounded diameter
\cite[Sections~3--6]{ChodoshLi2024}. The two-dimensional differential
inequality underlying our closed-surface diameter estimate is
equivalent to theirs \cite[Lemma~16]{ChodoshLi2024}. Under a pointwise
positive scalar-curvature lower bound, Liokumovich and Maximo constructed
singular foliations of compact three-manifolds with simultaneous area,
diameter, and genus control \cite{LiokumovichMaximo2023}.
Liokumovich and Wang extended area and diameter control to complete
noncompact three-manifolds \cite{LiokumovichWang2026}. These are stronger
geometric conclusions in their pointwise setting, and they also imply
graph-width bounds. Their level-set statements concern connected
components of fibers of proper Morse functions; passing to their Reeb
graphs distinguishes such components. The graph target in
Theorem~\ref{thm:main} likewise permits a uniform bound on whole fibers.

There are also sharp spectral band inequalities. For example,
Hirsch, Kazaras, Khuri, and Zhang obtain quantitative bounds for bands
from positivity of $-\Delta+c\Scal$, including $c=1/4$, under the
relevant orientability and topological hypotheses
\cite{HirschKazarasKhuriZhang2024}. Such an end-to-end distance estimate
is a different conclusion from a graph-width bound. We use the
weighted surface argument directly, keeping the spectral condition,
the topology of separating surfaces, and the final fiber estimate
within one proof.

\subsection{The proof and its useful components}

After scaling, assume $\lambda=1$. A positive supersolution of
$(-4\Delta_g+\Scal_g)v\ge v$ gives a smooth function $t=-2\log v$
such that
\begin{equation}\label{eq:intro-weighted-scalar}
 D(g,t):=\Scal_g+2\Delta_g t-|dt|_g^2\ge1.
\end{equation}
The function $t$ supplies the positive area density $e^{-t}$, without
changing the metric in which the final width is measured. All
minimization and smoothing arguments involving this density take
place in compact regions.

The first step is a weighted form of the stable-surface estimate.
If a closed connected two-sided surface is stable for weighted area
minus the weighted volume integral of a smooth pressure $\mu$, its
equation is $H-\nu t=\mu$, where $H$ is its mean curvature and $\nu$
is the chosen unit normal. A positive function for its stability operator converts
the ambient inequality into a two-dimensional inequality of the same
form as \eqref{eq:intro-weighted-scalar}. A conformal geodesic index
argument then gives an intrinsic diameter bound in the \emph{original}
surface metric whenever
$D(g,t)+\mu^2-2|d\mu|_g$ has a positive lower bound. The local version,
proved on relatively compact surface domains in
Section~\ref{sec:weighted-surfaces}, also prevents stable minimizing
sheets from escaping a sufficiently large fixed neighborhood.

We next cut along a locally finite family of closed stable weighted
minimal surfaces. The family is chosen so that no finite subfamily
disconnects $M$, and is maximal with this property. Every component
$P$ of the cut manifold then has the following separation property:
each compact connected two-sided surface meeting $\partial P$
transversely, with its boundary exactly on $\partial P$, separates
$P$. Section~\ref{sec:cutting} proves this by a localized minimizing
argument. On the orientation cover, antisymmetric integral currents
retain coorientations downstairs. A local area-nonincreasing
retraction at a cut boundary permits full-neighborhood comparisons,
so ordinary interior regularity applies even at boundary contacts.

The separation property allows small surfaces to control large
distance slabs. To construct those surfaces without a boundary
regularity problem, Section~\ref{sec:doubling-bubbles} doubles $P$ and
smooths the metric and density compatibly with reflection. The useful
smoothing variable is the logarithm of the weighted area density of
the normal slices. Weighted minimality makes its first normal
derivative vanish at the doubling locus, and the curvature inequality
survives smoothing. A reflection-invariant weighted $\mu$-bubble then
restricts to a finite collection of compact separating surfaces in $P$,
each with uniformly bounded diameter. The graph with complementary
regions as vertices and interface components as edges is a tree,
because each interface separates. The bubble places the basepoint and
a fixed distant slab component in different regions, so one interface
meets every path between them. Section~\ref{sec:slabs} uses this common
interface to bound the slab component's diameter.

Finally, overlapping slabs give a cover of multiplicity at most two,
and a subordinate partition of unity maps $P$ to a graph.
Section~\ref{sec:graph-gluing} joins these maps across the original
cuts, following the bounded-interface gluing principle of
\cite[Section~3.10.1, Lemma~(A)]{GromovFourLectures2023}.
Each cut surface is a sphere or a projective plane, so its map
to a finite graph is nullhomotopic. The homotopies are kept inside
finite subgraphs whose vertices correspond to cover sets meeting
the relevant collar. This extra condition ensures that all collar
pieces mapping to a given old graph point lie near one common cover
set. New edges join only the two collars belonging to a single cut.
These observations control entire fibers after gluing, rather than
only their connected components.

The weighted surface estimate, the compatible smoothing of doubles,
and the collar construction each isolate a useful part of the
argument. Their proofs require local elliptic and geometric measure
theory, but no uniform bounds on the geometry at infinity.
